\documentclass[border=14pt]{standalone}
\usepackage{circuitikz}
\usetikzlibrary{calc}
\ctikzset{bipoles/length=1.5cm}
\begin{document}
\begin{circuitikz}[american, line width=0.9pt]
  \coordinate (A) at (0,3.8);
  \coordinate (B) at (6.4,4.7);
  \coordinate (C) at (7.6,0.7);
  \coordinate (D) at (1.0,-0.3);
  % lazo cerrado irregular, elementos variados en diagonal (estilo libro)
  \draw (A) to[L=$u_1$] (B);
  \draw (B) to[R=$u_2$] (C);
  \draw (C) to[C=$u_3$] (D);
  \draw (D) to[V=$u_4$] (A);
  % munones: ramas que continuan al resto de la red
  \draw (A) -- ++(150:1.5) node[ocirc]{};
  \draw (A) -- ++(176:1.5) node[ocirc]{};
  \draw (B) -- ++(44:1.5) node[ocirc]{};
  \draw (B) -- ++(14:1.5) node[ocirc]{};
  \draw (C) -- ++(-30:1.5) node[ocirc]{};
  \draw (C) -- ++(-2:1.5) node[ocirc]{};
  \draw (D) -- ++(-108:1.5) node[ocirc]{};
  \draw (D) -- ++(-150:1.5) node[ocirc]{};
  % nodos
  \foreach \P in {A,B,C,D} \fill (\P) circle (3.2pt);
  \node[above left=2pt] at (A){$A$}; \node[above right=2pt] at (B){$B$};
  \node[right=3pt] at (C){$C$}; \node[below=3pt] at (D){$D$};
  % recorrido
  \draw[-{Latex[length=3.2mm]}, line width=1pt] (3.9,2.2) ++(110:1.05) arc (110:-205:1.05);
  \node[align=center] at (3.9,2.2){sentido de\\recorrido};
  \node[gray, align=center] at (9.9,3.4){\itshape resto\\\itshape de la red};
\end{circuitikz}
\end{document}
